\documentclass[12pt,a4paper]{article}

\usepackage[utf8]{inputenc}
\usepackage[T1]{fontenc}
\usepackage[french]{babel}
\usepackage{lmodern}
\usepackage[a4paper,margin=2.5cm]{geometry}
\usepackage{hyperref}
\usepackage{xcolor}
\usepackage{array}
\usepackage{booktabs}
\usepackage{tabularx}
\usepackage{enumitem}
\usepackage{amsmath,amssymb}
\usepackage{tcolorbox}
\usepackage{float}
\usepackage{parskip}
\usepackage{tikz}
\usetikzlibrary{arrows.meta,calc}

\definecolor{accent}{HTML}{2A5A8A}
\definecolor{muted}{HTML}{555555}
\definecolor{bgcode}{HTML}{F4F4F4}
\definecolor{warn}{HTML}{B85C00}
\definecolor{thbg}{HTML}{EEF3F8}

\hypersetup{colorlinks=true, linkcolor=accent, urlcolor=accent, citecolor=accent}

\tcbuselibrary{skins,breakable}
\newtcolorbox{remarque}{
  colback=thbg, colframe=accent, boxrule=0.8pt,
  left=8pt, right=8pt, top=6pt, bottom=6pt,
  title={\textbf{Remarque}}, fonttitle=\small
}

\title{\textbf{Implémentation de RTS-LLM}}
\author{Samir El Khoumri, Aly Hachem-Reda, Sami Chbicheb}
\date{Mai 2026}

\begin{document}

\maketitle
\tableofcontents
\newpage

% ============================================================
\section{Implémentation}
% ============================================================

L'implémentation de l'approche s'est déroulée en deux étapes itératives.
Nous avons d'abord développé une version reposant uniquement sur un modèle
de langage, afin de valider le principe de base. Nous avons ensuite enrichi
cette première version d'une phase d'analyse statique en amont, pour pallier
les limites observées.

% ============================================================
\subsection{Première version : sélection par le LLM seul}
% ============================================================

La première implémentation suit une logique directe : à partir d'un commit,
extraire le diff Git et l'ensemble des scénarios BDD du projet, puis
soumettre ces deux éléments à un LLM en lui demandant d'identifier les
scénarios susceptibles d'être impactés.

\subsubsection{Fonctionnement général}

Pour chaque commit analysé, le pipeline se déroule en trois étapes :

\begin{enumerate}[noitemsep]
  \item \textbf{Extraction du diff.} La commande \texttt{git diff} est
        exécutée entre le commit parent et le commit courant. Le diff
        obtenu contient les lignes ajoutées et supprimées dans chaque
        fichier modifié.

  \item \textbf{Collecte des scénarios BDD.} L'outil parcourt les fichiers
        \texttt{.feature} du projet et en extrait l'ensemble des scénarios
        avec leurs steps Gherkin.

  \item \textbf{Appel au LLM.} Un prompt est construit en concaténant le
        diff et la liste des scénarios numérotés. Le modèle est invité à
        retourner les indices des scénarios dont le comportement pourrait
        être affecté par les modifications, accompagnés d'une justification.
\end{enumerate}

\subsubsection{Construction du prompt}

Le prompt soumis au LLM positionne le modèle comme expert en tests de
régression BDD, et lui présente deux blocs d'information : le diff du
commit et la liste complète des scénarios. Il lui est demandé de répondre
au format JSON structuré :

\begin{center}
\ttfamily\small
\{ "selected": [2, 5], "reasoning": "..." \}
\end{center}

Ce format contraint la réponse et facilite son exploitation automatique.
La réponse est ensuite parsée pour en extraire les indices sélectionnés,
qui sont validés contre la liste des scénarios fournis afin d'éviter
toute hallucination d'indice.

\subsubsection{Modèle utilisé et configuration}

L'outil a été conçu pour fonctionner avec plusieurs fournisseurs de LLM
(OpenAI, Anthropic, Gemini, Ollama). Pour l'évaluation, nous avons
utilisé \textbf{Mistral 7B} via Ollama, qui permet une exécution locale
sans dépendance à une API payante. La température est fixée à $0{,}1$
pour garantir un comportement quasi-déterministe et améliorer la
reproductibilité des résultats.

\subsubsection{Limites observées}

Cette première approche présente deux limites importantes, mises en
évidence lors de l'évaluation :

\begin{itemize}
  \item \textbf{Passage à l'échelle difficile.} Le prompt contient
        l'intégralité des scénarios du projet. À mesure que ce nombre
        croît, le contexte soumis au modèle s'allonge, ce qui dégrade
        la qualité du raisonnement et augmente les coûts d'inférence.

  \item \textbf{Manque de précision.} Sans connaissance de la structure
        du code, le LLM ne peut relier un changement dans une méthode
        à un scénario qu'en s'appuyant sur des indices lexicaux dans le
        diff. Il tend à sur-sélectionner (faux positifs) par prudence,
        ou à manquer des impacts indirects non apparents dans le diff
        (faux négatifs).
\end{itemize}

Ces limites ont motivé le développement d'une seconde version intégrant
une phase d'analyse statique en amont.

% ============================================================
\subsection{Seconde version : analyse statique en amont du LLM}
% ============================================================

La seconde version introduit une couche d'analyse statique du code source
avant l'appel au LLM. L'idée centrale est de ne soumettre au modèle que
les scénarios déjà présélectionnés sur la base d'un lien structurel
démontrable avec le code modifié, réduisant ainsi à la fois le bruit et
la taille du contexte.

\subsubsection{Analyse statique du projet}

Avant tout appel au LLM, l'outil analyse le code source du projet cible
à l'aide de la bibliothèque \textbf{JavaParser}. Cette analyse produit
deux structures :

\begin{itemize}
  \item Un \textbf{graphe d'appels} reliant les méthodes entre elles :
        un arc de $A$ vers $B$ signifie que la méthode $A$ appelle la
        méthode $B$. Ce graphe est également construit en sens inverse
        (\emph{reverse call graph}) pour permettre de remonter de
        n'importe quelle méthode vers ses appelants.

  \item Une \textbf{table de correspondance} entre les scénarios BDD et
        les méthodes Java de leurs step definitions : pour chaque scénario,
        on connaît l'ensemble des méthodes Java invoquées par ses étapes
        Gherkin.
\end{itemize}

\subsubsection{Identification des scénarios candidats}

Lorsqu'un commit est soumis, les méthodes modifiées par le diff sont
d'abord identifiées par croisement des lignes modifiées avec les plages
de méthodes issues de l'analyse statique. On obtient ainsi un ensemble
$M$ de méthodes touchées par le commit.

À partir de $M$, on remonte le reverse call graph par un parcours en
largeur (\emph{BFS}) pour identifier tous les composants transitivement
dépendants des méthodes modifiées. Un scénario est retenu comme candidat
si au moins une de ses step definitions appartient à cet ensemble. On
obtient ainsi un sous-ensemble $C \subseteq \mathcal{S}$ des scénarios
structurellement liés aux modifications.

\begin{remarque}
  Par construction, tout scénario qui invoque --- directement ou
  indirectement --- une méthode modifiée est inclus dans $C$. La couche
  statique est donc \emph{sûre} : elle ne peut pas exclure un scénario
  véritablement impacté, sous réserve que le graphe d'appels soit
  suffisamment complet.
\end{remarque}

\subsubsection{Raffinement par le LLM}

Le LLM reçoit ensuite uniquement les scénarios de $C$, et non plus
l'ensemble du projet. Pour chaque candidat, le prompt inclut également
la \textbf{chaîne d'appel} qui relie sa step definition à la méthode
modifiée, par exemple :

\begin{center}
\ttfamily\small
setBodyTo $\to$ AbstractBddStepDefinition.addQueryParameters
\end{center}

Cette information permet au LLM de comprendre \emph{comment} le scénario
atteint le code modifié, et non seulement \emph{que} ce lien existe. Il
peut ainsi distinguer un impact réel d'un faux positif structurel (lien
existant dans le graphe, mais modification sans effet sur le comportement
observable du scénario, comme un changement de log ou de commentaire).

\subsubsection{Apport de la seconde version}

\begin{figure}[H]
\centering
\begin{tikzpicture}[
  >=Stealth, thick, font=\small,
  box/.style={
    draw=accent, fill=accent!8, rounded corners=4pt,
    minimum width=3.8cm, minimum height=1cm, align=center
  },
  arr/.style={->, thick}
]

\node[box] (diff)   at (0, 0)    {Diff Git};
\node[box] (static) at (4.5, 0)  {Analyse statique\\$\rightarrow C \subseteq \mathcal{S}$};
\node[box] (llm)    at (9, 0)    {Raisonnement LLM\\$\rightarrow S \subseteq C$};

\draw[arr] (diff)   -- node[above,font=\footnotesize]{méthodes modifiées} (static);
\draw[arr] (static) -- node[above,font=\footnotesize]{candidats + chaînes} (llm);

\draw[->, dashed, color=muted, thick]
  (static.south) -- ++(0,-0.9)
  node[below,font=\footnotesize\itshape]{fallback \texttt{--no-llm}}
  -- ++(2.5,0) -- ++(0,0.9);

\end{tikzpicture}
\caption{Pipeline de la seconde version : filtrage statique puis raffinement LLM.}
\label{fig:pipeline-v2}
\end{figure}

La combinaison des deux couches apporte trois améliorations concrètes
par rapport à la version LLM seul :

\begin{itemize}[noitemsep]
  \item Le contexte soumis au LLM est réduit aux seuls scénarios
        pertinents, ce qui améliore la qualité du raisonnement et
        réduit la latence.
  \item Le LLM dispose d'une information structurelle précise (chaînes
        d'appel) qu'il ne pouvait pas déduire seul à partir du diff.
  \item En cas d'indisponibilité du LLM, l'analyse statique seule
        (option \texttt{--no-llm}) constitue un repli valide qui préserve
        la sûreté de la sélection.
\end{itemize}

\end{document}
